\section{Appendix: Scope and Qualifications}

This concise study concerns the universal-calendar Resumed Sixth Sunday after the Epiphany, 15 November 2026, in the 1962 Roman Missal. The second-class Sunday has Gloria, Creed and the Trinity Preface, without an ordinary commemoration of Saint Albert the Great. Local and institutional additions are outside its scope. The complete appointed texts, their historical English witnesses and the fuller source discussion belong to the expansive study. The Epistle is 1 Thessalonians 1:2--10; the erroneous citation printed with the resumed Mass is corrected by its actual body and the corresponding winter lesson.

The two whole-Mass interpretations are editorial syntheses from the accepted study. Attributions concern each author's identified passage. Augustine's full-psalm exposition extends beyond the verses appointed as chants; Bede explains the surrounding faith and forgiveness sequence. The bounded research established no direct patristic commentary on the exact three orations or Trinity Preface. Schuster supplies the stated later commentary on the prayers. The source limits and textual checks are recorded in the canonical research and appointed-text audit.

Historical English is identified by edition and is not an approved English translation of the 1962 Missal. Public-domain historical wording retains its own United States source basis; no general permission in modern deliveries or packaging is asserted. Aquinas's argument is explained without reproducing the restricted modern editorial expression. This is a study aid, not an official liturgical book or an ecclesiastically approved commentary.

\section{References}

\begin{itemize}
\item \textbf{Appointed texts.} \textit{Missale Romanum}, Vatican typical edition, 1962, pp.~421--422, nos.~1759--1768; Trinity Preface, p.~293, no.~1082; corresponding Epistle, p.~49, no.~338. \url{https://media.churchmusicassociation.org/pdf/missale62.pdf}. Biblical contexts follow the Douay--Rheims (Challoner), Project Gutenberg no.~1581. Psalm numbering is Vulgate with Hebrew in parentheses. The historical English expression used to explain the Collect is from Eugene Cummiskey, publisher, \textit{The Roman Missal Translated into the English Language for the Use of the Laity}, Philadelphia, 1861, Sixth Sunday after Epiphany, p.~102.

\item \textbf{John Chrysostom.} \textit{Homilies on Matthew}, 46.2--3 and 47.1, historical English in \textit{Nicene and Post-Nicene Fathers}, first series, vol.~10, ed. Philip Schaff, CCEL text. \textit{Homilies on First Thessalonians}, homily 1, NPNF first series, vol.~13; its preface attributes the Thessalonian translation to James Tweed.

\item \textbf{Jerome.} \textit{Commentary on Matthew}, II, on 13:31--35, \textit{Patrologia Latina} 26 (Paris, 1845), cols.~89--93. \textit{Commentary on Jeremiah}, V, on 29:10--14, PL 24 (Paris, 1865), cols.~893--894. Original English exposition of the identified Latin.

\item \textbf{Augustine.} \textit{Expositions of the Psalms}, 43.6--8, 84.2--6 and 129.1--8, historical English in NPNF first series, vol.~8, ed. Philip Schaff, CCEL text. The edition heads these Psalms 44, 85 and 130. The morning-watch explanation follows Augustine's own lemma and exposition.

\item \textbf{Thomas Aquinas.} \textit{Super I Epistolam ad Thessalonicenses lectura}, chapter 1, lecture 1 [87862], Turin 1953 Latin edition, Corpus Thomisticum delivery revised by Enrique Alarcón, \url{https://www.corpusthomisticum.org/c1t.html}. Used through attributed exposition.

\item \textbf{Bede.} \textit{In Marci evangelium expositio}, on Mark 11:22--25, Migne-derived Latin Wikisource witness, paragraphs beginning \textit{Et respondens Jesus} and \textit{Et cum stabitis ad orandum}; contextual exposition of faith, charity, pardon and eternal life.

\item \textbf{Robert Bellarmine.} \textit{Commentary on the Book of Psalms}, trans. John O'Sullivan, Dublin: James Duffy, 1866, eCatholic2000 historical-text transcription; Psalm 43, explanations 7--9; Psalm 84, explanation 1; Psalm 129, explanations 3--5. Explanation numbers follow the commentary. The translation omits some Latin philological discussion.

\item \textbf{Ildefonso Schuster.} \textit{The Sacramentary}, trans. Arthur Levelis-Marke, London: Burns Oates \& Washbourne, I (1924), pp.~416--417, on the Collect, Secret and Postcommunion; III (1927), p.~188, on the Offertory \textit{De profundis}. The historical English is distinguished from the Italian original.
\end{itemize}
